\section{Outcome (d): Mutually exclusive events}

\subsection{What the outcome asks}

The exam tests whether you can recognise when two or more events cannot occur together
(mutually exclusive, or disjoint, events) and whether you then add their probabilities
correctly. It also tests the reverse direction: given $\Prob(A)$, $\Prob(B)$ and
$\Prob(A \cup B)$, decide whether the events are disjoint, or solve for whatever
intersection probability makes the numbers consistent. Underneath both directions sits one
technique that drives almost every Topic 1 problem: split the event you want into
non-overlapping pieces, find the probability of each piece, and add. Partitions of the
sample space (a family of disjoint events that together cover everything) are the standard
tool for that split. The phrases ``exactly one'', ``at least one'', ``at most one'' and
``none'' are all evaluated by writing the event as a union of disjoint pieces and adding.

\subsection{Definitions and rules}

Two events $A$ and $B$ are \emph{mutually exclusive} (equivalently \emph{disjoint}) if
$A \cap B = \varnothing$, so they cannot both occur on the same trial. A family
$A_1, A_2, \dots$ is \emph{pairwise disjoint} if $A_i \cap A_j = \varnothing$ for every
$i \neq j$. A family is \emph{exhaustive} if its union is the whole sample space $S$. A
family that is both pairwise disjoint and exhaustive is a \emph{partition} of $S$: every
outcome lies in exactly one of the events.

\begin{keyfact}[Addition rule for disjoint events]
If $A \cap B = \varnothing$ then
\[
\Prob(A \cup B) = \Prob(A) + \Prob(B).
\]
More generally, if $A_1, \dots, A_n$ are pairwise disjoint then
\[
\Prob\!\left(\bigcup_{i=1}^{n} A_i\right) = \sum_{i=1}^{n} \Prob(A_i),
\]
and the same holds for a countably infinite family $A_1, A_2, \dots$ (countable
additivity, one of the probability axioms). Additivity is the \emph{only} rule that lets you
add probabilities directly; it applies exactly when the events are disjoint.
\end{keyfact}

\begin{keyfact}[Complement as a partition]
For any event $A$, the pair $\{A, \cc{A}\}$ is a partition of $S$, so
\[
\Prob(A) + \Prob(\cc{A}) = 1, \qquad \Prob(\cc{A}) = 1 - \Prob(A).
\]
More generally, if $B_1, \dots, B_n$ partition $S$ then
$\Prob(B_1) + \dots + \Prob(B_n) = 1$. This is how the exam hides one probability: it gives
$n-1$ of them and expects you to recover the last one by subtraction.
\end{keyfact}

\begin{keyfact}[Disjoint decomposition identities]
For any events $A$ and $B$ (disjoint or not), the following unions are unions of disjoint
pieces, so their probabilities add:
\begin{align*}
A &= (A \cap B) \cup (A \cap \cc{B}),
 & \Prob(A) &= \Prob(A \cap B) + \Prob(A \cap \cc{B}), \\
A \cup B &= A \cup (B \cap \cc{A}),
 & \Prob(A \cup B) &= \Prob(A) + \Prob(B \cap \cc{A}), \\
A \cup B &= (A \cap \cc{B}) \cup (A \cap B) \cup (\cc{A} \cap B),
 & \Prob(A \cup B) &= \Prob(A \cap \cc{B}) + \Prob(A \cap B) + \Prob(\cc{A} \cap B).
\end{align*}
The first line is the law of total probability in its simplest form: condition on whether
$B$ happens or not. With a partition $B_1, \dots, B_n$ it reads
$\Prob(A) = \sum_{i} \Prob(A \cap B_i)$.
\end{keyfact}

\begin{keyfact}[Inclusion--exclusion collapses when events are disjoint]
For any two events, $\Prob(A \cup B) = \Prob(A) + \Prob(B) - \Prob(A \cap B)$. For three,
\[
\Prob(A \cup B \cup C) = \Prob(A) + \Prob(B) + \Prob(C)
 - \Prob(A \cap B) - \Prob(A \cap C) - \Prob(B \cap C) + \Prob(A \cap B \cap C).
\]
When the events are pairwise disjoint every intersection has probability $0$ and the
formula reduces to plain addition. Conversely, if the given numbers satisfy
$\Prob(A \cup B) = \Prob(A) + \Prob(B)$ then $\Prob(A \cap B) = 0$; if
$\Prob(A \cup B) < \Prob(A) + \Prob(B)$ the events overlap with positive probability.
\end{keyfact}

\begin{keyfact}[Disjoint is not independent]
If $A$ and $B$ are mutually exclusive and $\Prob(A) > 0$, $\Prob(B) > 0$, then $A$ and $B$
are \emph{dependent}.

\emph{Proof.} Disjointness gives $\Prob(A \cap B) = \Prob(\varnothing) = 0$, while
$\Prob(A)\Prob(B) > 0$. Hence $\Prob(A \cap B) \neq \Prob(A)\Prob(B)$, which is the
definition of dependence. Equivalently, $\Prob(A \mid B) = 0 \neq \Prob(A)$: knowing $B$
occurred rules $A$ out completely. \qed
\end{keyfact}

The three adjectives the exam uses for a family of events answer three different questions.

\begin{center}
\begin{tabular}{@{}lll@{}}
\toprule
Term & Definition & Consequence for probabilities \\
\midrule
Mutually exclusive & $A \cap B = \varnothing$ & $\Prob(A \cap B) = 0$; $\Prob(A \cup B) = \Prob(A) + \Prob(B)$ \\
Independent & $\Prob(A \cap B) = \Prob(A)\Prob(B)$ & $\Prob(A \mid B) = \Prob(A)$; can both occur \\
Exhaustive & $A \cup B = S$ & $\Prob(A \cup B) = 1$; at least one always occurs \\
\bottomrule
\end{tabular}
\end{center}

Mutually exclusive is a statement about the \emph{sets} (they do not overlap). Independent
is a statement about the \emph{numbers} (the product rule holds). Exhaustive is a statement
about \emph{coverage} (nothing is left out). Any combination is possible except mutually
exclusive together with independent when both probabilities are positive.

\subsection{Disjoint decomposition as a technique}

The standard move on a Topic 1 word problem is to rewrite the event asked for as a union of
pieces that cannot overlap, then add the probabilities of the pieces. Two things make the
move reliable. First, the pieces are usually intersections of the given events and their
complements, so each piece is a single ``cell'' of the Venn diagram. Second, every cell
probability can be recovered from the quantities the exam typically gives: the marginals
$\Prob(A)$, $\Prob(B)$, $\Prob(C)$, the pairwise intersections and the triple intersection.

\paragraph{Two events.} The four cells $A \cap B$, $A \cap \cc{B}$, $\cc{A} \cap B$ and
$\cc{A} \cap \cc{B}$ have probabilities $\Prob(AB)$, $\Prob(A) - \Prob(AB)$,
$\Prob(B) - \Prob(AB)$ and $1 - \Prob(A) - \Prob(B) + \Prob(AB)$, where $\Prob(AB)$
abbreviates $\Prob(A \cap B)$. ``Exactly one'' is $(A \cap \cc{B}) \cup (\cc{A} \cap B)$,
two disjoint cells, so its probability is $\Prob(A) + \Prob(B) - 2\Prob(AB)$. ``At least
one'' adds the cell $A \cap B$ back in; ``at most one'' removes it from the whole space;
``neither'' is the single cell $\cc{A} \cap \cc{B}$.

\paragraph{Three events.} Let $S_1 = \Prob(A) + \Prob(B) + \Prob(C)$,
$S_2 = \Prob(AB) + \Prob(AC) + \Prob(BC)$ and $S_3 = \Prob(ABC)$. The eight cells of the
three-set Venn diagram are disjoint. ``Exactly two'' is the union of the three cells
$A \cap B \cap \cc{C}$, $A \cap \cc{B} \cap C$ and $\cc{A} \cap B \cap C$; each pairwise
intersection contains the triple intersection, so each cell has probability (pairwise)
$- S_3$, and adding gives $S_2 - 3S_3$. Cell-by-cell reasoning of the same kind gives the
rest of the list below.

\begin{keyfact}[``Exactly'', ``at least'', ``at most'' formulas]
Two events, with $\Prob(A \cap B)$ written $\Prob(AB)$:
\begin{align*}
\Prob(\text{exactly one}) &= \Prob(A) + \Prob(B) - 2\Prob(AB), \\
\Prob(\text{at least one}) &= \Prob(A) + \Prob(B) - \Prob(AB), \\
\Prob(\text{at most one}) &= 1 - \Prob(AB), \\
\Prob(\text{none}) &= 1 - \Prob(A) - \Prob(B) + \Prob(AB).
\end{align*}
Three events, with $S_1 = \Prob(A) + \Prob(B) + \Prob(C)$,
$S_2 = \Prob(AB) + \Prob(AC) + \Prob(BC)$ and $S_3 = \Prob(ABC)$:
\begin{align*}
\Prob(\text{exactly three}) &= S_3, \\
\Prob(\text{exactly two}) &= S_2 - 3S_3, \\
\Prob(\text{exactly one}) &= S_1 - 2S_2 + 3S_3, \\
\Prob(\text{at least one}) &= S_1 - S_2 + S_3, \\
\Prob(\text{at least two}) &= S_2 - 2S_3, \\
\Prob(\text{at most one}) &= 1 - S_2 + 2S_3, \\
\Prob(\text{none}) &= 1 - S_1 + S_2 - S_3.
\end{align*}
Check on any problem: exactly one $+$ exactly two $+$ exactly three $=$ at least one.
\end{keyfact}

When the given events are themselves mutually exclusive, $S_2 = S_3 = 0$ and every formula
above collapses: ``at least one'' and ``exactly one'' coincide and equal $S_1$, ``at most
one'' has probability $1$, and ``none'' has probability $1 - S_1$.

\subsection{Worked examples}

\begin{example}[Claim size bands]
An insurer classifies each claim into exactly one of four size bands. The probability that a
claim is small (under 1{,}000) is $0.42$, medium (from 1{,}000 to under 5{,}000) is $0.28$,
and large (from 5{,}000 to under 25{,}000) is $0.18$. All other claims are catastrophic.
Calculate the probability that a randomly selected claim is either medium or catastrophic.

\smallskip\noindent (A)~$0.40$\hfill (B)~$0.46$\hfill (C)~$0.58$\hfill (D)~$0.34$\hfill (E)~$0.30$\hfill\mbox{}
\end{example}
\begin{solution}
The four bands partition the sample space, so their probabilities sum to $1$:
\[
\Prob(\text{catastrophic}) = 1 - 0.42 - 0.28 - 0.18 = 0.12.
\]
Medium and catastrophic are disjoint bands, so the addition rule applies directly:
\[
\Prob(\text{medium} \cup \text{catastrophic}) = 0.28 + 0.12 = 0.40.
\]
Answer: \textbf{(A)}.
\end{solution}

\begin{example}[Are the events disjoint?]
For a randomly chosen customer of an insurer, the probability of holding an auto policy is
$0.35$, the probability of holding a homeowners policy is $0.45$, and the probability of
holding at least one of the two is $0.70$. Calculate the probability that the customer holds
both policies.

\smallskip\noindent (A)~$0.00$\hfill (B)~$0.20$\hfill (C)~$0.15$\hfill (D)~$0.10$\hfill (E)~$0.05$\hfill\mbox{}
\end{example}
\begin{solution}
Let $A$ be auto and $B$ be homeowners. If $A$ and $B$ were disjoint, $\Prob(A \cup B)$ would
equal $0.35 + 0.45 = 0.80$. The stated value $0.70$ is smaller, so the events overlap and
the general addition rule is required:
\[
\Prob(A \cap B) = \Prob(A) + \Prob(B) - \Prob(A \cup B) = 0.35 + 0.45 - 0.70 = 0.10.
\]
Answer: \textbf{(D)}. The gap between the disjoint-case value $0.80$ and the actual value
$0.70$ is exactly the intersection probability.
\end{solution}

\begin{example}[Exactly one of two]
A policyholder has both an auto policy and a homeowners policy with the same company. In a
given year the probability of a claim on the auto policy is $0.40$, the probability of a claim
on the homeowners policy is $0.25$, and the probability of claims on both is $0.10$.
Calculate the probability that the policyholder files a claim on exactly one of the two
policies.

\smallskip\noindent (A)~$0.65$\hfill (B)~$0.45$\hfill (C)~$0.60$\hfill (D)~$0.35$\hfill (E)~$0.55$\hfill\mbox{}
\end{example}
\begin{solution}
Let $A$ be an auto claim and $H$ a homeowners claim. ``Exactly one'' is the disjoint union
$(A \cap \cc{H}) \cup (\cc{A} \cap H)$. Each piece is a marginal minus the intersection:
\[
\Prob(A \cap \cc{H}) = 0.40 - 0.10 = 0.30, \qquad
\Prob(\cc{A} \cap H) = 0.25 - 0.10 = 0.15.
\]
Adding the disjoint pieces gives $0.30 + 0.15 = 0.45$, which matches the formula
$\Prob(A) + \Prob(H) - 2\Prob(A \cap H) = 0.40 + 0.25 - 0.20 = 0.45$.
Answer: \textbf{(B)}. Choice (E) is $\Prob(A \cup H) = 0.55$, the ``at least one'' answer.
\end{solution}

\begin{example}[Exactly two of three]
A group health plan offers hospitalisation ($A$), dental ($B$) and vision ($C$) riders. For a
randomly selected member, $\Prob(A \cap B) = 0.12$, $\Prob(A \cap C) = 0.09$,
$\Prob(B \cap C) = 0.07$ and $\Prob(A \cap B \cap C) = 0.03$. Calculate the probability that
the member holds exactly two of the three riders.

\smallskip\noindent (A)~$0.25$\hfill (B)~$0.28$\hfill (C)~$0.22$\hfill (D)~$0.16$\hfill (E)~$0.19$\hfill\mbox{}
\end{example}
\begin{solution}
``Exactly two'' is the disjoint union of three cells, each a pairwise intersection with the
triple intersection removed:
\begin{align*}
\Prob(A \cap B \cap \cc{C}) &= 0.12 - 0.03 = 0.09, \\
\Prob(A \cap \cc{B} \cap C) &= 0.09 - 0.03 = 0.06, \\
\Prob(\cc{A} \cap B \cap C) &= 0.07 - 0.03 = 0.04.
\end{align*}
Sum: $0.09 + 0.06 + 0.04 = 0.19$. Equivalently $S_2 - 3S_3 = 0.28 - 0.09 = 0.19$.
Answer: \textbf{(E)}. Choice (B) is $S_2$ itself, which triple-counts the members holding
all three riders.
\end{solution}

\begin{example}[Mutually exclusive with an unknown]
An applicant for life insurance is classified as either preferred or standard, never both,
and some applicants are declined. Let $A$ be the event that an applicant is classified
preferred and $B$ the event that an applicant is classified standard. The probability that an
applicant is accepted (preferred or standard) is $0.60$, and an applicant is twice as likely
to be classified preferred as standard. Calculate $\Prob(B)$.

\smallskip\noindent (A)~$0.40$\hfill (B)~$0.30$\hfill (C)~$0.20$\hfill (D)~$0.15$\hfill (E)~$0.60$\hfill\mbox{}
\end{example}
\begin{solution}
The problem states that $A$ and $B$ cannot both occur, so they are mutually exclusive and
$\Prob(A \cup B) = \Prob(A) + \Prob(B)$. With $\Prob(A) = 2\Prob(B)$,
\[
0.60 = 2\Prob(B) + \Prob(B) = 3\Prob(B) \quad\Longrightarrow\quad \Prob(B) = 0.20,
\]
and $\Prob(A) = 0.40$. Answer: \textbf{(C)}. The probability of being declined is
$1 - 0.60 = 0.40$, which is choice (A); read the question stem carefully.
\end{solution}

\subsection{Traps}

\begin{trap}[Adding without disjointness]
$\Prob(A \cup B) = \Prob(A) + \Prob(B)$ is valid only when $A \cap B = \varnothing$. If the
problem does not state that the events cannot occur together (``exactly one of'', ``never
both'', ``classified into one of the following categories''), use the general rule
$\Prob(A) + \Prob(B) - \Prob(A \cap B)$. A quick sanity check: if the plain sum exceeds $1$,
the events cannot be disjoint. If the sum is at most $1$, the events might still overlap;
only the problem statement or a given intersection settles it.
\end{trap}

\begin{trap}[Disjoint is not independent]
Students frequently treat ``mutually exclusive'' as a synonym for ``independent''. It is
close to the opposite. Disjoint events with positive probability are always dependent,
because $\Prob(A \cap B) = 0$ while $\Prob(A)\Prob(B) > 0$. When a question says ``mutually
exclusive'', set every intersection to $0$; when it says ``independent'', set every
intersection to the product of the marginals. Never do both.
\end{trap}

\begin{trap}[Conditioning on a disjoint event]
If $A$ and $B$ are mutually exclusive and $\Prob(B) > 0$ then
$\Prob(A \mid B) = \Prob(A \cap B)/\Prob(B) = 0$. Likewise
$\Prob(A \cup \cc{B}) = \Prob(\cc{B})$ because $A \subseteq \cc{B}$. Answer choices in such
questions are built from the numbers $\Prob(A)$ and $\Prob(B)$ to catch anyone who computes
$\Prob(A)/\Prob(B)$ or $\Prob(A) + \Prob(\cc{B})$ mechanically.
\end{trap}

\begin{trap}[``Exactly'' versus ``at least'']
``Exactly one'' excludes the intersection; ``at least one'' includes it. For two events the
difference is exactly $\Prob(A \cap B)$, and the exam usually lists both values among the
choices. For three events, ``exactly two'' is $S_2 - 3S_3$, not $S_2$: every pairwise
intersection already contains the triple intersection, so $S_2$ counts the triple cell three
times. ``At least two'' is $S_2 - 2S_3$, which keeps the triple cell once.
\end{trap}

\subsection{Practice problems}

\begin{practice}
Claims on a commercial property policy fall into exactly one of four bands. The probability a
claim is under 500 is $0.35$, from 500 to under 2{,}000 is $0.40$, and 10{,}000 or more is
$0.05$. Calculate the probability that a claim is at least 2{,}000.

\smallskip\noindent (A)~$0.45$\hfill (B)~$0.20$\hfill (C)~$0.05$\hfill (D)~$0.25$\hfill (E)~$0.30$\hfill\mbox{}
\end{practice}

\begin{practice}
For a randomly selected policyholder, the probability of filing a medical claim in a year is
$0.30$, the probability of filing a prescription-drug claim is $0.50$, and the probability of
filing at least one of the two is $0.65$. Calculate the probability that the policyholder
files exactly one of the two types of claim.

\smallskip\noindent (A)~$0.50$\hfill (B)~$0.35$\hfill (C)~$0.65$\hfill (D)~$0.45$\hfill (E)~$0.55$\hfill\mbox{}
\end{practice}

\begin{practice}
Each policy in a portfolio is in exactly one of three risk classes: $A$, $B$ or $C$, or is
unclassified. Let $A$, $B$ and $C$ also denote the events that a randomly chosen policy is in
the respective class. You are given $\Prob(A) = 3\Prob(C)$, $\Prob(B) = 2\Prob(C)$ and
$\Prob(A \cup B \cup C) = 0.90$. Calculate $\Prob(A \cup C)$.

\smallskip\noindent (A)~$0.55$\hfill (B)~$0.45$\hfill (C)~$0.70$\hfill (D)~$0.75$\hfill (E)~$0.60$\hfill\mbox{}
\end{practice}

\begin{practice}
A survey of policyholders finds that $50\%$ carry collision coverage ($A$), $40\%$ carry
comprehensive coverage ($B$) and $30\%$ carry rental reimbursement ($C$). Also,
$20\%$ carry both $A$ and $B$, $15\%$ carry both $A$ and $C$, $10\%$ carry both $B$ and $C$,
and $5\%$ carry all three. Calculate the probability that a randomly selected policyholder
carries exactly two of the three coverages.

\smallskip\noindent (A)~$0.45$\hfill (B)~$0.30$\hfill (C)~$0.35$\hfill (D)~$0.25$\hfill (E)~$0.40$\hfill\mbox{}
\end{practice}

\begin{practice}
For a randomly selected insured, let $A$ be the event of a claim on a first-party coverage
and $B$ the event of a claim on a third-party coverage. You are given
$\Prob(A \cap B) = 0.12$, $\Prob(A \cap \cc{B}) = 0.28$ and $\Prob(\cc{A} \cap B) = 0.18$.
Calculate the probability that the insured has no claim on either coverage.

\smallskip\noindent (A)~$0.60$\hfill (B)~$0.30$\hfill (C)~$0.40$\hfill (D)~$0.42$\hfill (E)~$0.58$\hfill\mbox{}
\end{practice}

\begin{practice}
Events $A$ and $B$ are mutually exclusive with $\Prob(A) = 0.30$ and $\Prob(B) = 0.50$.
Calculate $\Prob(A \cup \cc{B})$.

\smallskip\noindent (A)~$0.50$\hfill (B)~$0.80$\hfill (C)~$0.30$\hfill (D)~$0.20$\hfill (E)~$0.65$\hfill\mbox{}
\end{practice}

\subsection{Answers to practice problems}

\begin{enumerate}[label=\arabic*.,nosep]
\item \textbf{(D)} The bands partition the sample space, so the missing band (2{,}000 to under
      10{,}000) has probability $1 - 0.35 - 0.40 - 0.05 = 0.20$. At least 2{,}000 is the union
      of the two disjoint top bands: $0.20 + 0.05 = 0.25$.
\item \textbf{(A)} $\Prob(A \cap B) = 0.30 + 0.50 - 0.65 = 0.15$. Exactly one is
      $\Prob(A \cup B) - \Prob(A \cap B) = 0.65 - 0.15 = 0.50$, or
      $0.30 + 0.50 - 2(0.15) = 0.50$.
\item \textbf{(E)} The classes are disjoint, so $3c + 2c + c = 6c = 0.90$ and $c = 0.15$.
      Then $\Prob(A \cup C) = 3c + c = 4c = 0.60$.
\item \textbf{(B)} $S_2 = 0.20 + 0.15 + 0.10 = 0.45$, $S_3 = 0.05$, and exactly two is
      $S_2 - 3S_3 = 0.45 - 0.15 = 0.30$. Check: exactly one is
      $1.20 - 0.90 + 0.15 = 0.45$, and $0.45 + 0.30 + 0.05 = 0.80 = S_1 - S_2 + S_3$.
\item \textbf{(D)} The three given cells are disjoint and make up $A \cup B$, so
      $\Prob(A \cup B) = 0.12 + 0.28 + 0.18 = 0.58$ and the answer is $1 - 0.58 = 0.42$.
      (Here $\Prob(A) = 0.40$, $\Prob(B) = 0.30$ and $\Prob(A)\Prob(B) = 0.12 = \Prob(A \cap B)$,
      so these events happen to be independent, and they are not disjoint.)
\item \textbf{(A)} Since $A \cap B = \varnothing$, $A \subseteq \cc{B}$, so
      $A \cup \cc{B} = \cc{B}$ and the probability is $1 - 0.50 = 0.50$. Choice (B) is
      the wrong result $\Prob(A) + \Prob(\cc{B}) = 0.80$, which double counts $A$.
\end{enumerate}

\subsection{One-page summary}

\begin{itemize}[nosep]
\item Mutually exclusive (disjoint): $A \cap B = \varnothing$, so $\Prob(A \cap B) = 0$ and
      $\Prob(A \mid B) = 0$. Pairwise disjoint: no two of the family overlap. Partition:
      pairwise disjoint and exhaustive; probabilities sum to $1$.
\item Addition rule: for disjoint events, $\Prob(A \cup B) = \Prob(A) + \Prob(B)$, and
      $\Prob(\bigcup A_i) = \sum \Prob(A_i)$ for any finite or countable pairwise disjoint
      family. This is the only situation in which probabilities add directly.
\item General rule: $\Prob(A \cup B) = \Prob(A) + \Prob(B) - \Prob(A \cap B)$. If the given
      union equals the sum of the marginals, the events are disjoint; if it is smaller, the
      shortfall is $\Prob(A \cap B)$.
\item Complement: $\Prob(\cc{A}) = 1 - \Prob(A)$. A partition with one probability missing:
      subtract the others from $1$.
\item Decomposition identities: $\Prob(A) = \Prob(A \cap B) + \Prob(A \cap \cc{B})$;
      $\Prob(A \cup B) = \Prob(A) + \Prob(\cc{A} \cap B)$;
      $\Prob(A \cup B) = \Prob(A \cap \cc{B}) + \Prob(A \cap B) + \Prob(\cc{A} \cap B)$.
\item Disjoint with positive probabilities implies dependent: $\Prob(A \cap B) = 0 \neq
      \Prob(A)\Prob(B)$. Mutually exclusive, independent and exhaustive are three unrelated
      properties.
\item Two events: exactly one $= \Prob(A) + \Prob(B) - 2\Prob(AB)$; at least one
      $= \Prob(A) + \Prob(B) - \Prob(AB)$; at most one $= 1 - \Prob(AB)$; none
      $= 1 - \Prob(A) - \Prob(B) + \Prob(AB)$.
\item Three events with $S_1, S_2, S_3$: exactly three $= S_3$; exactly two $= S_2 - 3S_3$;
      exactly one $= S_1 - 2S_2 + 3S_3$; at least one $= S_1 - S_2 + S_3$; at least two
      $= S_2 - 2S_3$; none $= 1 - S_1 + S_2 - S_3$.
\item Technique: rewrite the target event as a union of Venn-diagram cells, compute each cell
      from marginals and intersections, then add. Check that exactly one $+$ exactly two $+$
      exactly three $=$ at least one.
\item If $A$ and $B$ are disjoint then $A \subseteq \cc{B}$, so $\Prob(A \cup \cc{B}) =
      \Prob(\cc{B})$ and $\Prob(A \cap \cc{B}) = \Prob(A)$.
\end{itemize}
